\documentclass[11pt,a4paper]{article}
\usepackage[T1]{fontenc}
\usepackage[utf8]{inputenc}
\usepackage{lmodern}
\usepackage{amsmath,amssymb}
\usepackage[margin=25mm]{geometry}
\usepackage[hidelinks]{hyperref}
\hypersetup{pdftitle={Superdeterminism lives below every resolution; a floor gives it a lifetime},pdfauthor={navstokgap research working notes}}
\newcommand{\tightlist}{\setlength{\itemsep}{0pt}\setlength{\parskip}{0pt}}
\setlength{\emergencystretch}{3em}
\setcounter{secnumdepth}{0}
\begin{document}
\hypertarget{superdeterminism-lives-below-every-resolution-a-floor-gives-it-a-lifetime}{%
\section{Superdeterminism lives below every resolution; a floor gives it
a
lifetime}\label{superdeterminism-lives-below-every-resolution-a-floor-gives-it-a-lifetime}}

\textbf{Result, 2026-09-29 (Claude; written derivation, not yet
refereed).} Classical mechanics is deterministic, so the settings of a
Bell test are functions of the initial data; the statistical
independence of settings and hidden variables that Bell's theorem needs
is an assumption of typicality on those data, taken with the Liouville
measure. Superdeterminism is the choice of atypical data. The
proposition below shows where such data must hide when the setting
device is chaotic: in structure of the initial ensemble finer than any
fixed phase-space resolution. For a preparation ensemble resolved at
scale \(\delta\), the covariance between a hidden variable and a setting
chosen \(n\) steps later obeys

\[|{\rm Cov}_\rho(f,g\circ T^n)|\le2C_T\,K\,\|f\|_{\rm Lip}\|g\|_{\rm Lip}\,\delta^{-(d+1)}\theta^n,\]

so the correlation is below \(\eta\) after

\[n_*(\delta)=\frac{(d+1)\ln(1/\delta)+\ln\bigl(2C_TK\|f\|_{\rm Lip}\|g\|_{\rm Lip}/\eta\bigr)}{\ln(1/\theta)}\]

steps. Classical mechanics puts no lower bound on \(\delta\), and
\(n_*\to\infty\): a conspiracy can be stored in ever finer structure for
as long as needed. If a floor \(h\) on phase-space area per canonical
pair constrains the physical state, then
\(\ln(1/\delta)=\frac12\ln(A_0/h)\) with \(A_0\) the action scale of the
device, and the conspiracy has a finite lifetime
\(n_*\approx\frac{d+1}{2\ln(1/\theta)}\ln(A_0/h)\), the logarithmic
structure of the Ehrenfest time of quantum chaos
(\href{https://doi.org/10.1016/0378-4371(78)90190-5}{Berman and
Zaslavsky 1978}; metadata). With the floor, measurement independence for
chaotic setting devices is a theorem after \(n_*\); Bell's inequalities
then exclude local deterministic models that respect the floor. The
limit \(h\to0\) that Newton takes is exactly what reopens the loophole.

The user's remark that prompted this (2026-09-28/29): classical
mechanics is accidentally superdeterministic too, and the programme
holds that it has a mistake somewhere. The link is that typicality rests
on the Liouville measure, whose unit is an action per degree of freedom,
the same object whose unlimited refinement the
\href{principia-fifth-postulate.md}{fifth-postulate note} identifies as
Newton's joint determinacy.

\hypertarget{setting-and-the-proposition}{%
\subsection{1. Setting and the
proposition}\label{setting-and-the-proposition}}

Let \(M\) be a compact \(d\)-dimensional manifold with normalized volume
\(\mu\), and \(T\) a \(C^2\) volume-preserving Anosov diffeomorphism,
the setting device. Such maps are mixing with exponential decay of
correlations for Hölder, hence Lipschitz, observables
(\href{https://doi.org/10.1007/BFb0081279}{Bowen 1975}; metadata): with
\(\|F\|_{\rm Lip}=\sup|F|+{\rm Lip}(F)\),

\[\Bigl|\int F\,(G\circ T^n)\,d\mu-\int F\,d\mu\int G\,d\mu\Bigr|\le C_T\|F\|_{\rm Lip}\|G\|_{\rm Lip}\,\theta^n,\qquad\theta<1.\]

A preparation ensemble is a probability density \(\rho\) on \(M\); the
hidden variable of the measured system is \(\lambda=f(x)\) and the
setting is \(s=g(T^nx)\), with \(f,g\) Lipschitz. The ensemble is
\textbf{resolved at scale \(\delta\)} if
\(\|\rho\|_{\rm Lip}\le K\delta^{-(d+1)}\); for example any probability
measure smoothed by a mollifier of width \(\delta\), whose density has
height of order \(\delta^{-d}\) and slope of order \(\delta^{-(d+1)}\).

\textbf{Proposition.} For an ensemble resolved at scale \(\delta\),
\(|{\rm Cov}_\rho(f,g\circ T^n)|\le2C_TK\|f\|_{\rm Lip}\|g\|_{\rm Lip}\delta^{-(d+1)}\theta^n\).

\emph{Proof.} Apply the decay estimate twice, with \(F=f\rho\) and with
\(F=\rho\), using \(\int\rho\,d\mu=1\):
\(\int f\rho\,(g\circ T^n)=\int f\rho\int g+E_1\) and
\(\int\rho\,(g\circ T^n)=\int g+E_2\). Then
\({\rm Cov}_\rho=E_1-E_2\int f\rho\), and
\(\|f\rho\|_{\rm Lip}\le\|f\|_{\rm Lip}\|\rho\|_{\rm Lip}\),
\(|\int f\rho|\le\|f\|_{\rm Lip}\) give the bound. \(\square\)

\hypertarget{what-it-says-and-about-whom}{%
\subsection{2. What it says, and about
whom}\label{what-it-says-and-about-whom}}

\begin{itemize}
\tightlist
\item
  \textbf{Classical mechanics.} Nothing in the laws fixes \(\delta\); an
  ensemble can be a point mass. Independence is therefore an assumption
  about the data, and the superdeterministic option is a choice of data
  with structure below every scale. This is the precise sense of
  ``accidentally superdeterministic''.
\item
  \textbf{A floor.} The proposition needs the floor to constrain the
  physical state (its phase-space density), beyond constraining records.
  Under that premise, correlations of a chaotic setting device with a
  hidden variable decay after a time logarithmic in \(A_0/h\), the
  Ehrenfest-type scale at which classical and quantum chaos part.
\item
  \textbf{'t Hooft's automaton} lies outside the proposition. Its ontic
  states are exact cells and its correlations are carried by exact
  labels, with no Lipschitz structure to smooth; a floor on phase-space
  density says nothing about them
  (\href{discrete-substrate-actors.md}{discrete-substrate note}).
\item
  \textbf{Palmer's invariant set postulate} places states on a fractal
  invariant set and counterfactual settings off it
  (\href{https://doi.org/10.1098/rspa.2009.0080}{Palmer 2009}; read at
  the level of its title and our memory of the framework). A fractal
  support is structure at all scales, which a floor in the proposition's
  sense smears out; the two premises are in direct conflict. The
  superdeterminist position is surveyed in
  \href{https://doi.org/10.3389/fphy.2020.00139}{Hossenfelder and Palmer
  (2020)} (metadata).
\item
  \textbf{Quantum mechanics} is untouched: its Bell violations come from
  entangled states, and it assumes independent settings, consistently
  with the proposition.
\end{itemize}

\hypertarget{consequence-for-state}{%
\subsection{3. Consequence for STATE}\label{consequence-for-state}}

A small addition to the Newton side: the floor turns Bell's measurement
independence into a theorem for chaotic setting devices after an
Ehrenfest-type time, and the continuum limit \(h\to0\) reopens the
loophole. It is a statement about what the floor buys, with the floor
assumed; the necessity question is unchanged. To be refereed.
\end{document}
